\documentclass{article}

%Librerías
\usepackage[utf8]{inputenc}
\usepackage[spanish,mexico]{babel}
\setlength{\textwidth}{18cm}
\setlength{\oddsidemargin}{-1cm}
\setlength{\headsep}{-1cm}
\setlength{\voffset}{0cm}
\setlength{\topmargin}{0cm}
\setlength{\headheight}{0cm}
\usepackage{tikz}
\usetikzlibrary{calc,arrows}
\usepackage{multicol}
\usepackage{lipsum} 
%\bibliographystyle{apacite}
\usepackage[spanish]{babel}
\selectlanguage{spanish}


\begin{document}

%%%%%% ENCABEZADO %%%%%%%%%%%%%%%%%%%%%%%%%%%%%%%%%%%%%%%
% Logo de la maestría
\colorbox{white!10!}{
    \begin{minipage}[t]{0.05\textwidth} %0.165 
       \begin{flushright}
        \includegraphics[width=2in]{logo UPS.png}
       \end{flushright}
    \end{minipage}
    \begin{minipage}[H]{0.62 \textwidth} %0.62
        \begin{center}
         
        \end{center}
     \end{minipage}
    \begin{minipage}[t]{0.05 \textwidth}
        \begin{flushleft}
        \hspace{10.25cm}
            \includegraphics[width=2in]{Posgrados.png}
        \end{flushleft}
    \end{minipage}
}

\begin{tikzpicture}
    \draw[thick] (-6.5,0)--(11.2,0);
\end{tikzpicture}
%%%%%%%%%%%%%%%%%%%%%%%%%%%%%%%%%%%%%%%%%%%%%%%%%%%%%%%%%
\vspace{0.1cm}
\begin{center}
{\large\textsc{ANTEPROYECTO DEL TRABAJO DE TITULACIÓN}} \\
\vspace{0.5cm}
{ \large \textbf{Nombres y apellidos 1}} \\ 
\vspace{0.25cm}
{ \large \textbf{Nombres y apellidos 2 }}
\end{center}
\vspace{0.1cm}

\section{Tema del Trabajo de Titulación:  }
\begin{center}
Implementar Arquitectura orientada a servicios y DevOps para modulo de Historial Crediticio de la Cooperativa de ahorro y Crédito Jard	ín Azuayo.
\end{center}
\section{Docente tutor propuesto:   }
\begin{center}
PhD. Gonzalo Pérez
\end{center}
\section{Antecedentes:}
En este apartado se especificará de forma clara cómo nace la idea, si existen trabajos previos y cualquier otra información que permita conocer aspectos previos al trabajo propuesto  \cite{cataldi2000metodologia}.  

 \lipsum[1-1]
       
\section{Justificación:}   
Explica de forma breve y concisa la problemática que se aborda y destaca la importancia del trabajo a desarrollar. 
\section{Objetivos:}
Presenta el objetivo general del trabajo y los objetivos específicos que se alcanzarán con el desarrollo del mismo. 
\subsection*{Objetivo General:}
Diseñar .....
\subsection*{Objetivos Específicos:}

\begin{itemize}
    \item Construir ....
    \item Diseñar ....
\end{itemize}

\section{Alcance:}
\section{Metodología:}

Describe de manera clara el método empleado para cumplir los objetivos a través de las actividades dentro de la investigación. 
\section{Cronograma de actividades:}
Detalla las actividades generales que se desarrollarán hasta la fecha en la que el estudiante presentará su Trabajo de Titulación. 

\section{Presupuesto: (Opcional)}

Este apartado incluirá los recursos e información financiera para el Trabajo de Titulación. 
\bibliographystyle{abbrv}
%apacite



\bibliography{sample}

\vspace{1cm}
Este apartado contiene el listado de referencias empleadas para la elaboración del documento, de preferencia en formato APA u otros formatos Preconocidos.


\end{document}
